\documentclass{css}
%\documentclass[english]{css}

\usepackage[dvipdfmx]{graphicx}
\usepackage{latexsym}
\usepackage{listings}
\usepackage{url}
\urlstyle{same}

\def\|{\verb|}

\newcommand{\cssyear}[0]{2026}
\newcommand{\cssname}[0]{CSS 2026}
\newcommand{\cssversion}[0]{2026/06/01}
\newcommand{\cssemail}[0]{css2026-office@iwsec.org}

\makeatletter
\newcommand\newblock{\hskip .11em\@plus.33em\@minus.07em}
\makeatother

\begin{document}

%% 本文が和文の場合，タイトル・著者名・著者所属・概要は，和文・英文共に必須．

\title{脆弱性検知条件の供給源の再検討 \\ ---NVD・VulnCheckの追従性の定量評価とLLMを用いた一次データソースの構造化---}
\etitle{Revisiting the Supply of Vulnerability Detection Conditions: Quantifying the Timeliness of NVD and VulnCheck, and LLM-assisted Structuring of Primary Data Sources}

\affiliate{FUTURE}{フューチャー株式会社\\
Future Corporation}

%% メールアドレスは省略可能だが，代表者のメールアドレスは必須．
%% 姓名の間は半角スペースを入れること．

\author{中岡 典弘}{Norihiro Nakaoka}{FUTURE}[n.nakaoka.46@future.co.jp]
\author{篠原 俊一}{Shunichi Shinohara}{FUTURE}[s.shinohara.yx@future.co.jp]

\begin{abstract}
我々はOSS脆弱性スキャナVulsと脆弱性DBを開発・運用している．
CVE報告件数の急増（2026年は210日間で44,809件）を受け，検知条件をどこから得るべきかを定量的に明らかにする．
まずNVDとVulnCheck NVD++の全件（2026年7月29日時点，取下げ除き353,704件）を調査し，NVDによるCPE適用条件の付与率が2023年以前の100\%から2024年80.5\%，2025年62.6\%，2026年56.1\%へ低下し，2025年公開分の37.3\%をNVD自身が解析対象外と明示していることを示した．
次にFortinetのCVRF/CSAFと突き合わせ，Fortinet採番CVEでも一致は63.3\%にとどまること，製品に組み込まれたOpenSSH等のOSS部品の脆弱性では両者ともFortinet製品を列挙しないことを明らかにした．
さらにGit履歴から付与遅延を測り，NVDは中位数7日・42.6\%が未付与（Ciscoでは16日・48.3\%），VulnCheckは1日・22.6\%であった．
そこで，一次データソースをLLMで検証可能・再現可能に構造化する方式を提案する．
\end{abstract}

%% キーワード (1--5単語) の記載は任意．

\begin{jkeyword}
脆弱性検知，脆弱性データベース，NVD，一次データソース，LLM
\end{jkeyword}

\begin{eabstract}
We develop Vuls, an open-source vulnerability scanner, and its database.
Given the sharp growth of reported CVEs (44,809 in the first 210 days of 2026), we ask where detection conditions should come from and answer with measurements.
Analysing all 353,704 non-rejected entries of NVD and VulnCheck NVD++ as of 2026-07-29, we show NVD's coverage of CPE applicability statements fell from 100\% up to 2023 to 80.5\% (2024), 62.6\% (2025) and 56.1\% (2026), and that NVD marked 37.3\% of 2025 entries Deferred.
Matching against Fortinet's CVRF/CSAF, we further show only 63.3\% agree even for Fortinet-assigned CVEs, and that for vulnerabilities in embedded OSS components such as OpenSSH, neither source lists a single Fortinet product.
From our Git history we measure assignment latency at a median of 7 days for NVD with 42.6\% never assigned (16 days, 48.3\% for Cisco).
We therefore propose a verifiable, reproducible LLM-assisted method for structuring primary sources.
\end{eabstract}

%% the following keyword part is optional and can be omitted.

\begin{ekeyword}
Vulnerability Detection, Vulnerability Database, NVD, Primary Data Source, Large Language Model
\end{ekeyword}

\maketitle

\section{はじめに}
\label{sec:intro}

我々はオープンソース脆弱性スキャナ Vuls\cite{vuls} と，その脆弱性データベース群を継続的に開発・運用している．
前稿\cite{nakaoka2025}では，一次データソース（ベンダやプロジェクトが自ら公開するセキュリティアドバイザリ）を可能な限り直接収集し，取得したままの \texttt{vuls-data-raw} 形式と共通スキーマの \texttt{vuls-data-extracted} 形式に分けてGitで履歴管理する「データハーベスト」方式を提案した\cite{vuls-data-update,vuls-data-db}．
本稿はその基盤の上で，「検知条件を誰から得るべきか」という問いを定量的に扱う．

脆弱性スキャナが判定に用いる情報を本稿では\textbf{検知条件}と呼ぶ．
検知条件とは最低限「どの製品の」「どのバージョン範囲が」影響を受けるかの組であり，これが機械可読な形で与えられなければスキャナはインベントリと突き合わせられない．
NVD\cite{nist-nvd}においてこの役割を担うのが，\texttt{configurations} に格納されるCPE適用条件である．
各条件は \texttt{criteria} にCPE 2.3形式の文字列を持ち，任意で \texttt{versionStartIncluding} / \texttt{versionEndExcluding} 等の範囲指定を伴う．
ただしバージョン部が \texttt{*} や \texttt{-} で範囲指定を持たない場合，その条件は製品全体を影響対象とみなすことしか表現できない．

近年，CVEの報告件数は急増している．
我々の観測では2026年は1月1日から7月29日までの210日間で44,809件（取下げを除く）が公開されており，年率換算で約7.8万件，2023年の28,817件の約2.7倍のペースである．
この増加が解析体制を圧迫していることは，NVDが \texttt{Deferred} というステータスの運用を始めたこと\cite{nvd-status}に表れている．これは遅延ではなくNVDがエンリッチメントを行わないと明示した状態である．
またこの遅延を受け，VulnCheckはNVDのデータに独自の解析・エンリッチメントを加えたVulnCheck NVD++\cite{vulncheck-nvd2}の提供を開始した．NVD由来の \texttt{configurations} を保持したまま独自の結果を \texttt{vcConfigurations} と \texttt{vcVulnerableCPEs} に持つため，両者の寄与を機械的に切り分けられる．

我々のような検知ツールの提供者にとってこれは，「これまで信頼して使えていたデータソースを今後も同じ前提で使えるのか」という運用上の問題である．
しかし我々の知る限り，CVE増加後のNVDおよび商用エンリッチメントについて，検知条件の供給という観点から定量的に整理した報告はない．

そこで本稿は，NVDとVulnCheck NVD++の全件（2026年7月29日時点，取下げを除く353,704件）についてCPE適用条件の付与率と粒度を測り，さらに一次データソースを基準とする4基準を定義してFortinetとCiscoで実測する（\ref{sec:supply}章）．
そのうえで，LLMの非決定性を日々のデータ経路から追い出したまま一次データソースを構造化する3つの適用パターンを整理して実装し（\ref{sec:llm}章），入力データソースのみを変えた4種のデータベースで計10環境を検知して，上記の差が検知結果に現れることを確認する（\ref{sec:eval}章）．
NVDのみでの未検知率はWindows Server 2022の1.3\%からサポート終了版FortiOSの68.8\%まで開き，\ref{sec:supply}章で測った付与率はこの差をほとんど予測しない．

\section{背景}

\subsection{検知条件とCPE適用条件}

NVDの各CVEエントリは \texttt{configurations} に0個以上のCPE適用文を持つ．
各適用文は \texttt{cpeMatch} の集合であり，\texttt{criteria} にCPE 2.3形式の文字列，任意で \texttt{versionStartIncluding} / \texttt{versionEndExcluding} 等の範囲指定を持つ．
スキャナはインベントリ中の製品とバージョンがこの条件を満たすかで検知を行う．

ただし\textbf{CPE適用条件が存在することとバージョンを特定できることは同じではない}．
バージョン部が \texttt{*}（全バージョン）や \texttt{-}（指定なし）で範囲指定を伴わない場合，その条件は製品全体を影響対象とみなすことしか表現できない．
この区別は4.3節で定量化する．

\subsection{NVDの解析体制の変化とVulnCheck NVD++}

NVDはCVEエントリの処理状況を \texttt{vulnStatus} として公開している\cite{nvd-status}．
公式定義によれば \texttt{Awaiting Analysis} は「エンリッチメント対象として印が付けられた」状態，\texttt{Deferred} は「現時点でエンリッチメント予定に入っていない．NVDが宣言する対象範囲外か，資源等の事情で優先度が付けられていない」状態である．
すなわち \texttt{Deferred} は遅延ではなく，\textbf{NVDがエンリッチメントを行わないと明示した状態}である．

VulnCheck NVD++（本稿では \texttt{nist-nvd2} 系列）\cite{vulncheck-nvd2}は，NVDのCVEデータに独自の解析結果を加えて提供する．
NVD由来の \texttt{configurations} を保持しつつ独自エンリッチメント結果を \texttt{vcConfigurations} および \texttt{vcVulnerableCPEs} に持つため，「NVDが付与していないがVulnCheckは付与している」件数を機械的に切り分けられる．

\subsection{CNAが提供するaffected情報}

NVDのJSON 2.0フィードは，CVEレコードのCNAコンテナ由来の製品情報を \texttt{affected} フィールドとして保持している．
これはNVD API 2.0スキーマの正式な項目（\texttt{cve\_affected\_1.0.json}）であり，\texttt{vendor}，\texttt{product}，\texttt{packageName}，\texttt{cpes}，\texttt{defaultStatus}，および \texttt{versions[]}（\texttt{version}，\texttt{status}，\texttt{versionType}，\texttt{lessThan}，\texttt{lessThanOrEqual} 等）を持つ．
情報源ごとに配列要素が分かれ，CNAが記述した内容がほぼそのまま流れてくる．

つまりNVDは，一次データソース由来の製品情報を保持しながら，それをCPE適用条件へ翻訳する工程で詰まっている．
ただしそのバージョン表記は自然言語であり（例：\texttt{"2.5.1.720 and earlier"}，\texttt{"N/A"}），そのままでは照合に使えない（4.4節）．

\subsection{前提とする我々の基盤}

前稿\cite{nakaoka2025}で提案した fetch / extract / db build の3段構成を本稿でも前提とする．
fetch段は一次データソースの構造を保った \texttt{vuls-data-raw} 形式を，extract段は共通スキーマ \texttt{vuls-data-extracted} 形式を出力し，いずれもGitで履歴管理される．
\texttt{vuls-data-extracted} はエコシステム，パッケージ（またはCPE），影響を受けるバージョン範囲（比較演算子とバージョン文字列の解釈方法を明示）を表現できる．
執筆時点でfetch対象は77系統のベンダ・プロジェクト，196種類のデータソースに及ぶ（前稿時点では141種類）．
本稿の提案は，この基盤に「HTMLのみで提供される一次データソース」を取り込むための拡張である．

\section{関連研究}

Dongら\cite{dong2019}は，NVDの標準化されたエントリと，CVE概要文および参照先の非構造な脆弱性レポートとの不整合を検出するシステムVIEMを提案した．
深層学習に基づく固有表現抽出と関係抽出により製品名とバージョンを抽出し（precision 0.941, recall 0.993），約20年分・78,296件のCVE IDと70,569件のレポートを対象とした調査で，NVDエントリと影響バージョンが完全に一致するのは59.82\%にすぎないことを示した．

本研究はDongらと問題意識を共有するが，3点で異なる．
第一に，VIEMはNVDにエントリが存在することを前提に不整合を測るのに対し，本稿はCVE増加後に顕在化した\textbf{検知条件がそもそも生成されない}という失敗モード（\texttt{Deferred} / \texttt{Awaiting Analysis}）を対象とする．
第二に，VIEMは専用モデルの学習と表記変化ごとの再学習を要するのに対し，本稿は抽出精度単体ではなく\textbf{継続運用における保守性・再現性}を設計目標に置く．
第三に，VIEMは検出までを扱うのに対し，本稿は不整合を解消したデータベースを構築し検知精度への影響まで評価する．

NVDのデータ品質については，Anwarら\cite{anwar2022}が公開日・ベンダ名・深刻度等の不整合と欠損を体系的に調査し，NguyenとMassacci\cite{nguyen2013}がGoogle Chromeを対象に影響バージョンデータの系統的な誤りを示している．
いずれもCVE増加以前の調査であり，本稿は同種の問題が2024年以降に「欠損」として急拡大したことを示す点で相補的である．

\section{検知条件の供給実態の定量調査}

\subsection{調査対象と方法}

調査対象は，我々が運用するデータ収集基盤が取得した次の2系統のスナップショットである．

\begin{description}
  \item[NVD] \mbox{} \\
    NVD JSON 2.0 データフィード（\texttt{nvdcve-2.0-*.json.gz}）由来，371,416件\cite{nvd-feed}
  \item[VulnCheck] \mbox{} \\
    VulnCheck NVD++ (\texttt{nist-nvd2}) API 由来，371,393件\cite{vulncheck-nvd2}
\end{description}

スナップショットは2026年7月29日時点（データセット中の \texttt{lastModified} の最大値）である．
各データソースは取得時のコミットハッシュで固定しており，付録\ref{app:repro}のとおり同一条件で再検証できる．
両者のCVE ID集合を比較すると，NVDのみに存在するものが23件，VulnCheckのみに存在するものが0件であった．
\textbf{VulnCheck NVD++ の収録IDはNVDのID空間を超えない}．
したがってVulnCheckは，5章で定義するエントリ欠落（C1）を原理的に改善しない．

以降の集計では \texttt{vulnStatus} が \texttt{Rejected} の17,712件を除外し，353,704件を対象とする．
取下げエントリは製品情報を持たないことが正常であり，含めると付与率を不当に低く見せるためである\footnote{実例として，Fortinet採番のCVEには2025年に「Rejected reason: Not used」として一括取下げされたものが866件含まれる．除外しない場合Fortinetの付与率は32.8\%と算出されるが，除外後は99.6\%である．}．

集計は，各エントリについて \texttt{configurations}（NVD）と \texttt{vcConfigurations}（VulnCheck独自）を走査し，\texttt{vulnerable: true} の \texttt{cpeMatch} の個数とバージョン指定の種別を数えることで行った．
測定スクリプトは付録\ref{app:repro}に従い再現できる．

\subsection{CPE適用条件の付与率の推移}

表\ref{tab:year}に公開年別の解析状況とCPE適用条件の付与率を示す．
「NVD付与率」は \texttt{configurations} に \texttt{vulnerable} な \texttt{cpeMatch} を1つ以上持つ割合，「VulnCheck付与率」は \texttt{configurations} または \texttt{vcConfigurations} のいずれかに持つ割合であり，「VC補完」はNVDが付与しておらずVulnCheckのみが付与している件数である．

\begin{table*}[htbp]
  \caption{公開年別のCPE適用条件の付与状況（取下げを除く）}
  \ecaption{CPE applicability coverage by publication year (rejected entries excluded).}
  \label{tab:year}
  \begin{center}
  \begin{tabular}{rrrrrrrr}
    \hline
    公開年 & CVE件数 & 解析完了\% & 解析待ち\% & Deferred\% & NVD付与率\% & VulnCheck付与率\% & VC補完件数（\%） \\
    \hline
    2020 & 18,322 & 100.0 & 0.0 & 0.0 & 100.0 & 100.0 & 0 (0.0) \\
    2021 & 20,149 & 100.0 & 0.0 & 0.0 & 100.0 & 100.0 & 0 (0.0) \\
    2022 & 25,074 & 100.0 & 0.0 & 0.0 & 100.0 & 100.0 & 0 (0.0) \\
    2023 & 28,817 & 100.0 & 0.0 & 0.0 & 100.0 & 100.0 & 0 (0.0) \\
    2024 & 39,957 & 80.5 & 0.0 & 19.5 & 80.5 & 88.3 & 3,083 (7.7) \\
    2025 & 48,162 & 62.6 & 0.0 & 37.3 & 62.6 & 76.1 & 6,492 (13.5) \\
    2026$^\dagger$ & 44,809 & 55.9 & 9.3 & 34.8 & 56.1 & 78.5 & 10,058 (22.4) \\
    \hline
  \end{tabular}
  \end{center}
  \footnotesize
  $\dagger$ 2026年は1月1日〜7月29日（210日間）の公開分．
  「解析完了」は \texttt{Analyzed} と \texttt{Modified} の合計，「解析待ち」は \texttt{Received}，\texttt{Awaiting Analysis}，\texttt{Undergoing Analysis} の合計．
\end{table*}

\paragraph{(1) 2023年公開分までは付与率100\%である．}
2016年から2023年までのいずれの年も，取下げを除く全エントリにCPE適用条件が付与されている．
「報告件数が激増する前のNVDは信頼できるのか」という問いに対し，少なくとも検知条件の\textbf{有無}という観点では，2023年公開分以前は完全に補われていると答えられる．
ただし記述内容が正しいかは別問題であり，先行研究\cite{dong2019,nguyen2013}が指摘するとおり一致率は別途評価が必要である（7章）．

\paragraph{(2) 2024年を境に急落する．}
付与率は2024年80.5\%，2025年62.6\%，2026年56.1\%と低下し，低下分はほぼそのまま \texttt{Deferred} の増加（19.5\%，37.3\%，34.8\%）で説明される．
付与率と解析完了率がほぼ一致することから，NVDにおけるCPE適用条件の付与はエンリッチメント完了とほぼ同義であることも確認できる．

\paragraph{(3) \texttt{Deferred} は待っても解消しない．}
その定義は「エンリッチメント予定に入っていない」であり，時間が経てば付与される保証はない．
2025年公開分は本稿執筆時点で公開から最短でも7か月，平均では約1年半が経過しているが，37.3\%が \texttt{Deferred} のままである．
すなわち2025年以降のNVDについては，\textbf{約4割弱は将来にわたって検知条件が供給されないと想定して設計すべき}である．

\paragraph{(4) VulnCheckは補完するが埋め切らない．}
独自エンリッチメントは付与率を2024年88.3\%，2025年76.1\%，2026年78.5\%へ引き上げ，補完件数は年を追って増加する（2026年は10,058件，22.4\%）．
商用エンリッチメントの寄与は無視できない大きさである．
一方で2025年以降は依然として2割強が両者ともに未付与である．

なお，VulnCheckは2024年11月の記事で，2024年公開分のCPE付与率を自社76.95\%に対しNVD 41.35\%と報告している\cite{vulncheck-cpe}．
同じ2024年公開分を本稿が2026年7月に観測すると，NVD 80.5\%・VulnCheck 88.3\%である．
集計定義が同一とは限らないため厳密な比較ではないが，\textbf{NVDは公開直後は大きく遅れるものの，1年半をかけて差の多くを埋める}（両者の差は35.6ポイントから7.8ポイントへ縮小）ことを示唆する．
一方で，そこから先は \texttt{Deferred} として残り続ける．
したがってVulnCheckの価値は，最終的な網羅性よりも公開直後の期間に集中すると解釈できる．

公開月別のコホート（付録\ref{app:cohort}）で見ると，公開から約2年経過した2024年9月公開分は解析完了88.5\%・\texttt{Deferred} 11.5\%，公開から約7か月の2025年12月公開分は解析完了50.4\%・\texttt{Deferred} 49.6\%，公開から2週間の2026年7月公開分は解析完了33.5\%・解析待ち30.4\%・\texttt{Deferred} 36.0\%であった．
時間経過とともに解析待ちは解消されるが，\texttt{Deferred} の割合は月によって25\%から50\%の幅で残る．

\paragraph{(5) 未付与エントリは長く滞留する．}
NVDがエンリッチメント対象と認識している \texttt{Awaiting Analysis} でも公開から中位数43日・90\%点124日が経過しており，\texttt{Deferred} の41,363件は中位数286日，うち32,215件が90日を超える（付録\ref{app:survey-detail}）．
\textbf{新規に公開されたCVEについては，公開後1〜4か月はNVDの検知条件が存在しない前提で設計する必要がある}．
これは，公開直後の脆弱性ほど攻撃に利用されやすいという性質と正面から衝突する．

\subsection{付与された条件の粒度とCNA別の偏り}

\textbf{CPE適用条件が存在することとバージョンを特定できることは同じではない}．
2024年以降に公開された132,928件のうち，バージョン範囲または特定バージョンで表現されているのはNVDで63.9\%，VulnCheckを加えて77.8\%にとどまる（付録\ref{app:survey-detail}）．
VulnCheckの \texttt{vcConfigurations} には両者が同一エントリ内に混在する例もある．
CVE-2024-6387（regreSSHion）では \texttt{cpe:2.3:a:openssh:openssh} に 8.5p1〜9.7p1 という適切な範囲がある一方，\texttt{cpe:2.3:a:gnu:grub2:-} のようにバージョン指定を持たない条件が並び，そのまま用いれば当該製品を持つ全ホストが検知対象となる．
\textbf{付与率の向上が検知精度の向上を意味するとは限らない}．

採番CNA別に見ると差は極端である（付録\ref{app:survey-detail}）．
Microsoft 99.7\%，Fortinet 99.6\%，Apple 99.4\%と付与率がほぼ維持される一方，Cisco 77.2\%（\texttt{Deferred} 15.0\%），\textbf{JPCERT/CC採番は39.1\%（\texttt{Deferred} 59.7\%）}でVulnCheckを加えても48.4\%にとどまる．
JPCERT/CCはJVNを通じて国内ベンダの脆弱性を採番する経路であり，\textbf{国内ベンダ製品の検知条件はNVDもVulnCheckもほぼ供給していない}．
NVDの \texttt{Deferred} は無作為ではなく領域に偏って適用されている．
またMicrosoftは付与率99.7\%に対しバージョン特定可能率が92.9\%と乖離する．
WindowsのCPEはバージョン部が \texttt{-} で表現され，パッチ適用状態は更新プログラム（KB）情報に依存するためである．

\subsection{未付与時に残る情報：CNA提供のaffected}

NVDがCPE適用条件を付与していない45,413件（2024年以降公開，取下げを除く）について代替情報の有無を調べた．
結果は，CNA提供の \texttt{affected} が存在するもの45,413件（100.0\%），VulnCheckの \texttt{vcConfigurations} が存在するもの19,607件（43.2\%），いずれも存在しないもの0件であった．

すなわち\textbf{製品情報そのものは常に存在し，欠けているのはCPEへの翻訳だけである}．
しかしその内容はそのまま照合に使えない．実例を示す．

\begin{description}
  \item[CVE-2025-20271] \mbox{} \\
    Cisco採番，2025-06-18公開，\texttt{Deferred}．\texttt{configurations} および \texttt{vcConfigurations} は空．\texttt{affected} は \texttt{\{"vendor": "Cisco", "product": "Cisco Meraki MX Firmware", "versions": [\{"version": "N/A", "status": "affected"\}]\}} のみで，バージョンが \texttt{N/A} である．一方Ciscoのアドバイザリ本体には修正済みバージョンが明記されている．
  \item[CVE-2023-39933] \mbox{} \\
    JPCERT/CC採番，2024-03-18公開，\texttt{Deferred}．\texttt{affected} は \texttt{\{"vendor": "A.K.I Software ", "product": "pmc.exe", "versions": [\{"version": "2.5.1.720 and earlier", "status": "affected"\}]\}}．バージョンが自然言語（\texttt{and earlier}）で記述され，ベンダ名には末尾空白が含まれる．
\end{description}

これらは，一次データソースへ遡り，自然言語のバージョン表記を機械可読な範囲式へ正規化する必要があることを示す．
6章の提案はこの要求に応える．

\section{一次データソースとの比較}

4章はNVD・VulnCheck内部の情報だけによる調査である．
本章では一次データソースを基準とする4つの評価基準を定義し，FortinetとCiscoを対象に実測する．

\subsection{評価基準}

\begin{description}
  \item[C1（エントリ欠落）] \mbox{} \\
    一次データソースに存在する脆弱性が，NVD・VulnCheckに対応するエントリを持たない．
    両者はCVE IDを主キーとするため，CVE ID未採番の脆弱性は構造的に収録され得ない．
    そこでC1は，ベンダ独自ID（Fortinetの \texttt{FG-IR-yy-nnn}，Ciscoの \texttt{cisco-sa-*}，ヤマハのFAQページ等）のみで公開されている\textbf{C1a（未採番による欠落）}と，CVE IDは記載されているがエントリが存在しないか内容が空である\textbf{C1b（採番済みだが未収録）}に分ける．
    4.1節のとおりVulnCheckの収録IDはNVDのID空間を超えないためC1を改善せず，C1aはCVE制度の設計に起因するため一次データソースを直接収集する以外に解決手段がない．
  \item[C2（製品欠落）] \mbox{} \\
    一次データソースがaffectedとする製品が，NVD・VulnCheck側のaffected製品として列挙されていない．
  \item[C3（バージョン不一致）] \mbox{} \\
    affectedまたはfixedのバージョン範囲が両者で異なる．
    (a) 相手側が狭い（未検知を生む），(b) 相手側が広い（過剰検知を生む）に分類する．
    前稿\cite{nakaoka2025}のCVE-2023-26207（FortiOS 7.0.0〜7.0.15がFortinetでは影響ありとされるがNVDでは列挙されない）は(a)の例である．
  \item[C4（付与遅延）] \mbox{} \\
    一次データソースでの公開日から，検知条件が付与されるまでの日数．
\end{description}

\subsection{対象データと方法}

一次データソースには，ベンダが機械可読形式で公開しているものを用いる．
FortinetはPSIRT AdvisoryのCVRFおよびCSAF\cite{vuls-data-raw-fortinet-cvrf,vuls-data-raw-fortinet-csaf}，Ciscoは openVuln API 由来のJSON\cite{vuls-data-raw-cisco-json}である．
我々はFortinet PSIRTのHTMLを変換したデータ\cite{vuls-data-raw-fortinet}も保持しているが，CVRF/CSAFの更新に追従できておらず，2018〜2023年公開分ではノードの51\%でバージョンが空になる（実際のアドバイザリには記載がある）ため，本稿では基準として採用しない．
この追従漏れ自体は，手書き抽出器が提供元の変更で静かに壊れるという6章の動機の実例である．

Fortinet CVRF は \texttt{product\_tree} に影響を受けるバージョンを個別に列挙し（\texttt{FortiOS-7.6.1}，\texttt{FortiOS-7.6.0}…），\texttt{product\_statuses} で \texttt{Known Affected} を与える．
修正版は持たないため，未検知は測れるが過剰検知は測れない．
CSAF は \texttt{known\_affected} を範囲・単一版・「X.Y all versions」の3形式で与え，加えて \texttt{known\_not\_affected} として修正版を明示するため，過剰検知も測れる．
製品情報を伴うCVRFは2022年以降に限られ（2022年155件，2023年158件，2024年147件，2025年80件，2026年76件），CSAFは173件である．

Ciscoは影響を受けるバージョンを \texttt{productNames} に「Cisco Adaptive Security Appliance (ASA) Software 9.20.1」のような自然文で列挙する．
285,161件中97.3\%は末尾にバージョンを伴うが，\textbf{CPEを持たない}ため，C2・C3の判定には製品名からCPEへの対応付けが別途必要である．
本稿ではCiscoについてC1とC4のみを測定し，C2・C3は6章の方式の実装後に行う．

C2・C3の判定には\textbf{境界値法}を用いる．
バージョン範囲を連続区間として包含関係で比較すると誤判定が生じる．
例えばCVE-2025-22254では，Fortinetが \texttt{<=6.4.15} と記述する範囲をNVDは \texttt{versionEndExcluding: 6.4.16} と記述しており，実在するバージョンの集合としては同一だが，連続区間では「NVD側が広い」と判定されてしまう．
そこで (CVE, 製品) ごとに，一次データソースが影響ありとする版を\textbf{検知されるべき点}，修正版を\textbf{検知されるべきでない点}として，各点で相手側の条件の真偽が一致するかを評価する．
前者で相手が偽なら未検知，後者で相手が真なら過剰検知である．
製品の突合はCPEの製品部から区切り文字を除いた文字列で行う．
なお一次データソース側のバージョン表現が解析できない (CVE, 製品) 組は比較から除外した．

さらに結果をCVEの採番元で層別する．
Fortinetが採番したCVEと，Fortinet製品に組み込まれた部品（OpenSSH等のOSS）に対して他者が採番したCVE（以下，\textbf{組込み部品のCVE}と呼ぶ．Fortinetは自社製品が影響を受けることをアドバイザリで公表する）では，NVDの振る舞いが大きく異なるためである．
採番元はNVDの \texttt{sourceIdentifier} で判定した．

C4はGit履歴から測る．
我々はスナップショットを1日2回取得しGitで履歴管理しており\cite{vuls-data-db}，履歴はNVDが2025年6月26日，VulnCheckが2025年8月6日以降利用できる．
各CVEのファイルを履歴に沿って辿り，\texttt{vulnerable} な \texttt{cpeMatch} が初めて現れたコミットの日付を求めた．
起点はアドバイザリの \texttt{initial\_release\_date} である．

\subsection{C1：エントリ欠落}

\textbf{Fortinet}：CVRFの1,119件のアドバイザリに1,328件のCVEが対応する．
このうち\textbf{37件のアドバイザリはCVE IDを持たない}（C1a，FG-IR-17-271やFG-IR-17-259など2019年以前が中心）．
CVE IDを持つもののうちNVD・VulnCheckに存在しないものは0件であった（C1b）．

\textbf{Cisco}：5,309件のアドバイザリに6,538件のCVEが対応し，\textbf{260件（4.9\%）のアドバイザリはCVE IDを持たない}（C1a，\texttt{cisco-sa-20170515} や \texttt{cisco-sa-20070105-csacs} など2017年以前が中心）．
NVDに存在しないものは8件であった（C1b）．

\subsection{C2・C3：製品欠落とバージョン不一致}

結果を表\ref{tab:match}に示す．

\begin{table*}[htbp]
  \caption{Fortinet一次データソースとの一致状況（境界値法による，単位\%．各セルは NVD / VulnCheck）}
  \ecaption{Agreement with the Fortinet primary source (boundary-value method; each cell is NVD / VulnCheck).}
  \label{tab:match}
  \begin{center}
  \begin{tabular}{llrrrrrrrr}
    \hline
    基準 & 採番元 & CVE & 組 & 一致 & C2 製品欠落 & C3a 未検知 & C3b 過剰検知 & 未検知率 & 過剰検知率 \\
    \hline
    CVRF & Fortinet & 656 & 1,135 & 63.3 / 75.0 & 18.7 / 11.3 & 18.1 / 13.7 & --- & 17.1 / 9.7 & --- \\
    CVRF & 他者 & 54 & 255 & 0.0 / 0.0 & \textbf{100.0 / 100.0} & 0.0 / 0.0 & --- & 100.0 / 100.0 & --- \\
    CSAF & Fortinet & 168 & 298 & 71.1 / \textbf{89.9} & 15.8 / \textbf{0.0} & 9.4 / 6.4 & 3.7 / 3.7 & 10.8 / \textbf{1.5} & 0.5 / 0.5 \\
    CSAF & 他者 & 6 & 60 & 0.0 / 0.0 & \textbf{100.0 / 100.0} & 0.0 / 0.0 & 0.0 / 0.0 & 59.2 / 59.2 & 0.0 / 0.0 \\
    \hline
  \end{tabular}
  \end{center}
  \footnotesize
  一致・C2・C3a・C3bは(CVE,製品)組に対する割合，未検知率・過剰検知率は検証点に対する割合．
  CVRFは修正版を持たないため過剰検知（C3b）を測定できない．CVRFは製品情報を伴う2022年以降を対象とする．年別の内訳は付録D に示す．
\end{table*}

\paragraph{(1) 組込み部品の脆弱性では，NVDもVulnCheckもFortinet製品を1件も列挙しない．}
他者が採番したCVEについて，C2は\textbf{100.0\%}である（CVRFで54 CVE・255組，CSAFで6 CVE・60組）．
C2の内訳を見ると，その99〜100\%は「当該CVEにCPEは付与されているが，Fortinet製品のCPEが1件も含まれない」ケースであった．
該当例はCVE-2024-6387（regreSSHion），CVE-2023-48795（Terrapin），CVE-2024-3596（RADIUS）など，OpenSSHやRADIUS実装といった，Fortinet製品に組み込まれた部品の脆弱性である．
Fortinetは自社製品が影響を受けることをアドバイザリで公表しているが，NVDのCPEには部品そのものの製品（\texttt{openssh:openssh} 等）しか並ばない．
\textbf{OSS部品を組み込んだ製品を検知しようとすると，この欠落がそのまま未検知になる}．
4章の付与率にはこの欠落は一切現れない．

\paragraph{(2) Fortinetが採番したCVEでも一致は63.3\%（CVRF）／71.1\%（CSAF）にとどまる．}
4.5節でFortinetのCPE付与率は99.6\%と最良の水準であったが，内容の一致はそこまで高くない．
C2（18.7\%／15.8\%）の内訳は98〜99\%が「NVDにFortinet製品のCPEはあるが，当該製品のCPEが無い」ケースであり，FortiSOAR，FortiExtender，FortiClient Windows/Mac/Linux のように，Fortinetが区別する製品をNVDが列挙していない，あるいは異なる粒度で表す場合である．
\textbf{付与率は検知可能性の上限にすぎない}ことが，一次データソースとの直接比較で裏づけられた．

\paragraph{(3) VulnCheckはFortinetについて大きく改善する．}
CSAF基準・Fortinet採番では，一致が71.1\%から\textbf{89.9\%}へ，C2が15.8\%から\textbf{0.0\%}へ，検証点単位の未検知率が10.8\%から\textbf{1.5\%}へ改善する．
一方で過剰検知率は0.5\%とNVDと同等であり，\textbf{本稿の測定範囲では「VulnCheckが過剰検知を増やす」という傾向は認められなかった}．
ただしこれは修正版（\texttt{known\_not\_affected}）を負例とした評価であり，4.4節で見たバージョン指定を持たない粗いCPEの影響は，7章の実機・合成インベントリでの評価で改めて測る必要がある．

\paragraph{(4) 組込み部品のCVEはVulnCheckでも改善しない．}
他者採番のCVEについてはVulnCheckもC2 100.0\%であり，商用エンリッチメントもこの欠落を埋めていない．
すなわちOSS部品を組み込んだ製品の検知条件は，一次データソースから取得する以外に手段がない．

\subsection{C4：付与遅延}

表\ref{tab:delay}に，公開から検知条件が初めて現れるまでの日数を示す．

\begin{table*}[htbp]
  \caption{検知条件が付与されるまでの日数（Git履歴による）}
  \ecaption{Days until a detection condition first appears, measured from our Git history.}
  \label{tab:delay}
  \begin{center}
  \begin{tabular}{llrrrrrrl}
    \hline
    対象 & 起点 & 付与済み & 中位数 & 75\% & 90\% & 最大 & 30日超 & 打切り（未付与） \\
    \hline
    NVD（全体） & CVE公開日 & 1,592 & 7 & 16 & 50 & 219 & 245 & 1,181 / 2,773（42.6\%） \\
    VulnCheck（全体） & CVE公開日 & 1,973 & 1 & 4 & 11 & 238 & 75 & 576 / 2,549（22.6\%） \\
    NVD（Fortinet） & アドバイザリ公開日 & 184 & 3 & 7 & 14 & 101 & 4 & 2 / 186（1.1\%） \\
    VulnCheck（Fortinet） & アドバイザリ公開日 & 163 & 1 & 1 & 2 & 9 & 0 & 1 / 164（0.6\%） \\
    NVD（Cisco） & アドバイザリ公開日 & 151 & 16 & 57 & 85 & 125 & 65 & 141 / 292（48.3\%） \\
    VulnCheck（Cisco） & アドバイザリ公開日 & 224 & 1 & 3 & 44 & 125 & 35 & 26 / 250（10.4\%） \\
    \hline
  \end{tabular}
  \end{center}
  \footnotesize
  日数の単位は日．全体はGit履歴開始以降に公開されたCVEから25件ごとに抽出した標本，Fortinet・Ciscoは該当する全件．
\end{table*}

全体では，NVDは中位数7日・90\%点50日を要し，42.6\%は本稿執筆時点でも未付与である．
VulnCheckは中位数1日・90\%点11日で，未付与は22.6\%にとどまる．
\textbf{VulnCheckの最大の価値は最終的な網羅性ではなく初動の速さにある}という4.2節の解釈が，遅延の測定でも裏づけられた．

ベンダ間の差は大きい．
Fortinetでは NVD でも中位数3日・打切り1.1\%であるのに対し，\textbf{Ciscoでは中位数16日・90\%点85日で，48.3\%が本稿執筆時点でも未付与}である．
CiscoはVulnCheckでも打切り10.4\%が残る．
4.5節でCiscoは付与率77.2\%・\texttt{Deferred} 15.0\%，Fortinetは付与率99.6\%であった．
\textbf{付与率の低いベンダでは，欠落だけでなく遅延も深刻である}ことが確認できた．

\subsection{他の一次データソースへの拡張}

以上の比較が可能だったのは，一次データソースが機械可読な形で提供されているためである．
同じ比較を広げるには，HTMLのみで提供されるものを機械可読な形に変換する作業が必要になる．
表\ref{tab:sources}に対象とするデータソースを示す．
NVDでの供給状況（4章）が良好なもの（Fortinet，Microsoft，Apple）と不良なもの（Cisco，OpenSSL，JPCERT/CC経由のヤマハ）の両方を含め，「NVDで足りている領域」と「足りていない領域」を対比できるよう選定した．

\begin{table*}[htbp]
  \caption{対象とする一次データソース}
  \ecaption{Primary data sources targeted in this work.}
  \label{tab:sources}
  \begin{center}
  \begin{tabular}{lll}
    \hline
    分類 / ベンダ・プロジェクト & 一次データソースの提供形態 & 4章で観測されたNVDでの供給状況 \\
    \hline
    ネットワーク / Fortinet & CVRF/CSAF，PSIRT Advisory（HTML） & 付与率99.6\%．一致は63.3\%（5.4節） \\
    ネットワーク / Cisco & openVuln API（JSON，CPEなし），CVRF/CSAF & 付与率77.2\%，\texttt{Deferred} 15.0\%，遅延中位数16日 \\
    ネットワーク / ヤマハ & RTproのセキュリティFAQ（機種別HTML），JVN経由 & JPCERT/CC採番で付与率39.1\%，\texttt{Deferred} 59.7\% \\
    OS / Microsoft Windows & MSRC（CVRF/API） & 付与率99.7\%だがバージョン特定可能率92.9\% \\
    OS / Apple macOS & セキュリティリリースノート（HTML） & 付与率99.4\%（良好） \\
    ソフトウェア / Apache httpd & Security Report（HTML） & Apache CNA全体で付与率92.5\% \\
    ソフトウェア / OpenSSL & Vulnerabilities一覧・Advisory（HTML） & 付与率79.2\%，\texttt{Deferred} 20.8\% \\
    \hline
  \end{tabular}
  \end{center}
\end{table*}

Fortinet，Cisco，Microsoftについては既に機械可読形式からの収集を実装済みである\cite{vuls-data-update}．
一方，ヤマハ，Apple macOS，Apache httpd，OpenSSLはHTMLのみで提供されており未実装である．
次章の提案はこれらを対象とする．

\section{一次データソースの構造化}
\label{sec:llm}

\subsection{課題と設計要件}

一次データソースを直接収集しようとすると，構造が不安定であること（HTMLはスキーマを持たず，提供元の改修でセレクタが予告なく変わる）と，記述が多様で値が欠けていることの2つが問題になる．
LLMはこの両者に適するが，そのままデータ経路に置くと前稿で確立した検知結果の再現性を損なう．
そこで次の4要件を課す．
\begin{list}{}{%
  \setlength{\leftmargin}{1zw}\setlength{\itemindent}{-1zw}%
  \setlength{\labelwidth}{0pt}\setlength{\labelsep}{0pt}%
  \setlength{\topsep}{2pt}\setlength{\itemsep}{0pt}%
  \setlength{\parsep}{0pt}\setlength{\partopsep}{0pt}}
\item \textbf{R1（決定性）}同一入力から常に同一出力が得られる．
\item \textbf{R2（検証可能性）}出力を機械的に検証でき，通らないものはDBに入らない．
\item \textbf{R3（保守コストの局所化）}構造変更への追従が全面的な書き換えを要しない．
\item \textbf{R4（監査可能性）}生成元を追跡できる．
\end{list}
共通する設計はLLMの非決定性を日々のデータ経路から追い出すことであり，LLMに何を出力させるかで\textbf{P1}（データを一度だけ生成）・\textbf{P2}（抽出器のコードを生成）・\textbf{P3}（入力ごとに生成）の3つに整理した．この順はLLMが日々のデータ経路に残る度合いの順である．LLMは機械可読形式の代替ではなく，それが無いか不完全な場合の補完手段である．

\subsection{P1：欠損データの生成}
\label{sec:p1}

構造はあるが値が欠けている場合，LLMにデータを出力させ，凍結資産としてコミットする．実行は一度きりで日々のfetchでは呼ばれない（R1）．
Fortinet CVRFは2022年より前の \texttt{product\_statuses} が空で，1,119件のうち465件が該当する（値は\ref{sec:method3}節の影響を受けない版と同じく注記の自然文にしかない）．CSAFは2024年以降しか存在せず，この期間は一次データソースから検知できない．
そこでFortinet自身が書いた4つの記述——アドバイザリの2種の注記，CNAとして公開するCVEレコード，旧来の手書きデータセット——を突き合わせさせ，補完表を生成した\cite{vdu-pr892}．
検証は出典間の相互照合による（R2）．バージョン単位の照合，注記中のバージョン字句の網羅検査，NVDとの突合により計118件の不一致を検出し，根拠を記録しながら個別に裁定した（R4）．
バージョン制約を持たない19行は製品全体に一致してしまうため個別に監査し，未監査の行や未知の製品名は抽出を失敗させる——黙って落とすことも黙って広げることもしない．
結果として465件中422件が検知条件を得た．

\subsection{P2：抽出器の生成}
\label{sec:p2}

構造が安定していれば，fetchとextractのコードをLLMに書かせ，日々の収集は生成済みコードのみで決定的に動作させられる（R1）．検証はゴールデンセットへの回帰テストによる（R2）．
\ref{sec:eval}章で用いるApple Security Releases\cite{vdu-pr914}とMicrosoftの更新プログラム記事\cite{vdu-pr919}がこれにあたる．
このパターンの作業量は，構造の解析ではなく例外の発見と確定が占める．
Appleでは2003年以前のアーカイブ形式や公開が停止された記事の扱いが大半で，合計1,327行を要した．
このような網羅的な走査はLLMが得意とする作業であり，人手は生成された差分の承認に縮小される（R3）．
CIは取得件数の対前回比を閾値で監視してDBへの反映を止める．変更を再びLLMに与えて抽出器を直す運用は設計にとどまり，現状は人手で修正している．

\subsection{P3：HTMLからの直接変換}
\label{sec:p3}

構造が弱く鍵も欠けている場合，新規入力ごとにLLMにデータを出力させる．ここだけはデータ経路にLLMが残るため，変換結果をGitで持つことで決定性を作り（R1），型と検証器で検査する（R2）．
本稿が調べた範囲で最も極端なのはOpenSSHであった．
第一に，項目の境界が位置にしか現れない．セキュリティ情報のページ\cite{openssh-sec}の55件のエントリは \texttt{<li>} の中で日付・影響範囲・表題・本文が \texttt{<br>} で区切られているだけで \texttt{class} 属性も見出しも持たず，\texttt{<br>} の数は2個から7個までばらつくため，何番目の断片が影響範囲かをセレクタで固定できない．
第二に，鍵となるCVE IDがほとんど振られていない．本文にCVE IDを含むエントリは6件にとどまる．OpenSSHはCNAではなく，我々が観測した134件のCVEはMITREやRed Hatが採番している．アドバイザリを書く主体とCVEレコードを書く主体が別であることが欠落の理由である．
その帰結が\ref{sec:c1c3}節のCVE-2024-6387（regreSSHion）である．ページはこれを \texttt{8.5p1 to 9.7p1 (inclusive)} と正確な範囲付きで記述しながらCVE IDを書いていない．
検知条件は一次データソースだけで組み立てられるが，それが何と呼ばれている脆弱性なのかは分からない．

55件のエントリをすべて \texttt{vuls-data-raw} 形式へ変換した\cite{vdu-pr928}．
検査は2段ある．extractのコードは各バージョン文字列がOpenSSHのリリース表記かと範囲の上下限が高々1つで空でないかを検べ，通らなければ抽出を失敗させる．一方，出力の版とCVE IDが原文に逐語で現れるか，変換元の \texttt{<li>} を連結すると元のリストが復元されるかは，手順書がモデルに課す自己検査であり機械的に強制していない．原文に無い値の混入を止める検査こそがR2を満たしていない．
CVE IDの紐付けは一次データソースだけでは6件（10.9\%）にとどまるが，MITREのCVEレコードを参照することで45件（81.8\%）まで上がった．
一次データソースは影響範囲を持ちながら索引を持たず，MITREは索引を持ちながら影響範囲の精度を欠くという相補性が，実測として確認された．
ただしこの紐付けは検知件数を1件も変えていない．
生成元の追跡（R4）は，fetchがページを \texttt{origin/} に逐語で保存し，記録ごとに変換元の \texttt{<li>} を持たせることで満たす．変換の陳腐化——ページが変わったのに \texttt{raw/} を作り直していない状態——はCIで検出する\cite{vdb-pr224}．決定性（R1）は，変換結果である \texttt{raw/} をGitで持ち，fetchがそれを保存することで満たす．ページが変わらない限りモデルは走らず，出力は取得し直しても変わらない．残る不確かさは，ページが変わったときの再変換が同じ出力になる保証がないことである．

\section{検知精度への影響の評価計画}

5章は検知条件そのものを比較した．
本章はその差が実際のスキャン結果にどう現れるかを測る計画である．結果は発表時に報告する．

同一の基盤\cite{vuls-data-update,vuls-data-db}を用い，入力データソースのみを変えた3種のデータベースを構築する．
\textbf{DB-nvd} はNVDの \texttt{configurations} のみ，\textbf{DB-vc} はVulnCheck NVD++ の \texttt{configurations} と \texttt{vcConfigurations}，\textbf{DB-primary} は表\ref{tab:sources}の一次データソース（6章の方式で構造化したものを含む）である．

インベントリは2種を用いる．
\textbf{合成インベントリ}は5.2節の境界値法をそのまま用い，一次データソースの影響版と修正版から「検知されるべき」「検知されるべきでない」のラベル付きバージョンを生成する．
5章はこれを検知条件の集合演算で評価したが，本章では実際にスキャナを動かし，データベース構築と照合実装を含む経路全体での未検知・過剰検知を測る．
\textbf{実環境インベントリ}はVulsのスキャン対象として整備しているコンテナ・仮想環境\cite{vuls-targets}であり，運用に近い条件での差分を確認する．

指標は5.4節と同じ未検知率・過剰検知率に加え，スキャン結果の件数差，および過剰検知のうちバージョン指定を持たない条件（\texttt{-}，\texttt{*}）に起因する割合を報告する．
さらに公開年で2023年以前 / 2024年 / 2025年 / 2026年に分割して同じ指標を算出し，「どの時点までのNVDを信頼してよいか」を検知精度の観点から述べ，併せて4.5節のCNA別の差が検知精度の差として現れるかを確認する．

\section{議論と限界}

\paragraph{一次データソース側のデータ品質が測定の前提になる．}
当初はFortinet PSIRTのHTML変換データ\cite{vuls-data-raw-fortinet}を基準に測定していたが，このデータはCVRF/CSAFの更新に追従できておらず，2018〜2023年公開分ではノードの51\%でバージョンが空であった（実際のアドバイザリには記載がある）．
本稿の測定はすべてCVRF/CSAFに基づく．
一次データソースとの比較では，一次データソース側の抽出が健全であること自体を検証しなければならない．

\paragraph{製品の突合は綴りに依存する．}
CPEの製品部から区切り文字を除いて突合しているため，略称の差（\texttt{FortiClient Enterprise Management Server} と \texttt{forticlientems} 等）は自動では吸収できない．
表\ref{tab:match}のC2のうち「NVDに別のFortinet製品のCPEはある」ケース（Fortinet採番で98〜99\%）には，この対応付けの失敗が含まれる可能性がある．

\paragraph{観測期間の制約．}
4章は2026年7月29日時点の1点観測である．
5.5節のC4はGit履歴から測ったが，履歴はNVDが2025年6月26日，VulnCheckが2025年8月6日以降しか存在せず，それ以前の公開分は測れない．
また2026年後半の公開分は観測期間が短く，遅延の長い側が打ち切られる．

\paragraph{CPE適用条件の有無は検知可能性と等価ではない．}
4.3節および表\ref{tab:cna}のMicrosoftの例のとおり，条件が存在してもバージョンを特定できない場合がある．
逆にCPE辞書に製品が登録されていない場合は条件があっても検知できない．
本稿の付与率は検知可能性の上限を与えるものと解釈すべきである．

\paragraph{7章の「正解」は一次データソースである．}
未検知率・過剰検知率はDB-primaryを正解とみなした値であり，一次データソース側の過大な記載は過剰検知として現れる．
またMicrosoftのみ検知方式が異なり（適用済みKB対CPE），他の2ベンダと同列には比較できない．
スキャン対象は8環境と限られ，Microsoftの実機以外は宣言したCPEのみを持つ疑似サーバである．
またCiscoの製品名からCPEへの対応付けはアドバイザリの31.9\%しか覆っておらず，7章のCiscoの結果はIOSとIOS XEに限られる．

\paragraph{供給源の信頼度による抑制の代償と限界．}
VulnCheckのCPEは機械的に生成されるため，Microsoft Edgeのようにうまく機能するものと，Windowsのように大量の偽陽性を生むものが混在する（7.2節(2)(4)）．
Vulsが「検証済みソースが製品を定義していれば補助的ソースの検知を捨てる」規則を置いているのはこのためであり，規則そのものは必要な妥協である．
本稿が測ったのはその代償と限界である．
代償として，macOS 10.13.6では真陽性781件が捨てられている．
限界として，Windows 11 23H2の偽陽性833件はこの規則をすり抜けた．NVDが置いた抽象CPE（\texttt{all\_windows}）とVulnCheckが展開した具体CPE（\texttt{windows\_11\_23h2}）が別製品として扱われるためである．
規則を緩めれば前者は救えるが後者は悪化する．
必要なのは信頼度の調整ではなく，照合不能な抽象CPEをソース側で明示的に扱うことだと考えられる．

\paragraph{LLM利用に伴うリスク．}
6章のP1・P2は非決定性を生成時に閉じ込めるが，未実装のP3はデータ経路にLLMが残るため，キャッシュと検証器に決定性を依存する．
いずれの検証器も形式的な整合性を確認できるが「原文の意味を正しく解釈したか」は保証しない．
P1で出典間の相互照合により118件の不一致を検出したことは，単一の出力を信じてはならないことの裏返しでもある．
人手レビューのキューを設計に含めており，そのレビュー負荷が運用可能な水準に収まるかは未検証である．

その他の限界は付録\ref{app:limits}に記す．

\section{おわりに}

本稿では，脆弱性検知に不可欠な検知条件が現在どこから供給されているかを定量的に調査した．
NVDのCPE付与率は2023年公開分までの100\%から2024年以降56.1\%まで低下し，低下分はNVDが明示的に解析対象外とした \texttt{Deferred} に対応する．
しかし付与率は検知精度を予測しない．
組込み部品の脆弱性についてNVDもVulnCheckもFortinet製品のCPEを1件も列挙せず，Fortinetが採番したCVEでも一致は63.3\%にとどまった．
入力データソースのみを変えた4種のデータベースで10環境を検知すると，NVDのみでの未検知率は製品によって1.3\%から68.8\%まで開き，VulnCheckは付与率を上げる一方でWindowsに70.6\%の過剰検知を生んだ．
誤り方も製品で違い，Ciscoでは未検知53.7\%と過剰検知23.2\%が同時に，Appleでは下限を持たないCPEによる過剰検知81.5\%が生じた．
誤りは2023年以前のCVEに限っても減らず，「CVE増加前のNVDなら信頼してよい」という仮説は成り立たない．
併せて，我々自身が置いている抑制規則の代償（Cisco IOS XEでは抑制112件中31件が真陽性）と限界（偽陽性833件のすり抜け）を定量化した．

一次データソースを構造化する手段としては，欠損データの生成・抽出器の生成・HTMLからの直接変換の3パターンを整理し，Fortinet・Apple・Microsoft・OpenSSHに実装した．
今後の課題は，一次データソースから取得すべき製品とNVDやVulnCheckで十分な製品を見分ける方法を，より広い製品で確立することである．取得元を製品ごとに判断できなければ，一次データソースの収集コストは正当化できない．


\bibliographystyle{junsrt}
\bibliography{./refs.bib}

%% 研究倫理に関する記載はページ数制限の対象外．付録の前に置く．
\section{研究倫理}
\label{sec:ethics}

本研究は，公開されている脆弱性情報のみを対象とし，脆弱性の新規発見や実システムへの攻撃を伴わない．
調査に用いたデータは，NVDが公開するJSON 2.0データフィード，VulnCheck NVD++のAPI，および各ベンダ・プロジェクトが公開するセキュリティアドバイザリである．
いずれも一般に公開された情報であり，個人情報や非公開情報は含まない．

VulnCheck NVD++のデータ利用にあたっては，提供元が定める利用条件に従い，当該データがVulnCheck NVD++に由来することを明示する．
また各ベンダの公開情報の取得にあたっては，提供元サーバに過度な負荷をかけないよう取得間隔を制御する．

6章で提案するLLMの利用においては，LLMに入力するのは公開されたアドバイザリの本文のみであり，非公開情報を入力しない．
また，LLMの出力を検証せずに脆弱性データベースへ投入しない設計としており，誤った脆弱性情報が利用者に届くリスクを低減している．

本稿で報告する集計値は，公開データに対する機械的な集計であり，特定のベンダや組織の評価・非難を意図するものではない．
NVDの \texttt{Deferred} は運用方針の表明として公開されているものであり，本稿はそれを所与の条件として利用者側の対応を論じている．


%% 付録は本文の後ろ，最大5ページ．
\appendix
\section{測定の再現手順}
\label{app:repro}

本稿4章の数値は，以下の手順で再現できる．

\begin{enumerate}
  \item \texttt{vuls-data-update}\cite{vuls-data-update} により，NVD JSON 2.0フィード（\texttt{nvd/feed/cve/v2}）とVulnCheck NVD++（\texttt{vulncheck/nist-nvd2}）を取得する．または \texttt{vuls-data-db}\cite{vuls-data-db} のコンテナイメージから同名のタグを取得する．
  \item 本稿と同時に公開する集計スクリプト（\texttt{paper/analysis/}）を実行する．
    \texttt{extract.py} は全エントリから \texttt{id}，\texttt{published}，\texttt{vulnStatus}，\texttt{sourceIdentifier}，\texttt{cpeMatch} 数などをCSVに出力する．
    \texttt{agg3.py} は表\ref{tab:year}，表\ref{tab:pending}，表\ref{tab:cna}および4.6節の値を，\texttt{extract2.py} は表\ref{tab:gran}の値を算出する．
  \item 5章については，\texttt{fortinet\_c123.py} が表\ref{tab:match}を，\texttt{cisco\_c14.py} がCiscoのC1を，\texttt{c4.py}・\texttt{fortinet\_c4.py}・\texttt{cisco\_c4.py} が表\ref{tab:delay}を算出する．
    \texttt{c4.py} 系は \texttt{git log -{}-name-only} で対象CVEを含むコミットを列挙し，\texttt{git cat-file -{}-batch} で各時点の内容を取得して，\texttt{vulnerable} な \texttt{cpeMatch} が初めて現れたコミット日を求める．
  \item スナップショットの基準日は，データセット中の \texttt{lastModified} の最大値（本稿では2026年7月29日）とする．
  \item 本稿の数値を厳密に再現するには，各データソースのコミットを固定する必要がある．
    使用した全データソースのタグ名とコミットハッシュはリポジトリの \texttt{README.md} に記録しており，\texttt{sh paper/analysis/dataset-manifest.sh} で照合できる．
    主要なものは NVD が \texttt{f34a9c3b16}，VulnCheck が \texttt{cc430fb117}，Fortinet CVRF が \texttt{41b12c7484}，Fortinet CSAF が \texttt{cedaf51587}，Cisco が \texttt{5fdc0f3ea0} である（いずれも2026年7月30日取得）．
\end{enumerate}

\section{公開月別コホートの解析状況}
\label{app:cohort}

\begin{table}[htbp]
  \caption{公開月別の解析状況（取下げを除く，2026年7月29日時点）}
  \ecaption{Enrichment status by publication month (rejected entries excluded).}
  \label{tab:cohort}
  \begin{center}
  \small
  \begin{tabular}{lrrrrrrr}
    \hline
    公開月 & 件数 & 経過日 & 完了\% & Def.\% & 待ち\% & NVD\% & VC\% \\
    \hline
    2024-07 & 3,113 & 744 & 83.2 & 16.8 & 0.0 & 83.2 & 90.7 \\
    2024-08 & 2,897 & 713 & 87.5 & 12.5 & 0.0 & 87.5 & 93.6 \\
    2024-09 & 2,514 & 682 & 88.5 & 11.5 & 0.0 & 88.5 & 94.6 \\
    2024-10 & 3,566 & 652 & 77.8 & 22.2 & 0.0 & 77.8 & 86.5 \\
    2024-11 & 4,056 & 621 & 70.0 & 30.0 & 0.0 & 70.0 & 78.9 \\
    2024-12 & 3,431 & 591 & 59.0 & 41.0 & 0.0 & 59.0 & 71.0 \\
    2025-01 & 4,275 & 560 & 59.3 & 40.7 & 0.0 & 59.3 & 69.8 \\
    2025-02 & 3,673 & 529 & 71.8 & 28.2 & 0.0 & 71.8 & 78.3 \\
    2025-03 & 4,007 & 501 & 64.3 & 35.7 & 0.0 & 64.3 & 73.9 \\
    2025-04 & 4,034 & 470 & 57.2 & 42.8 & 0.0 & 57.2 & 69.8 \\
    2025-05 & 3,980 & 440 & 70.7 & 29.3 & 0.0 & 70.7 & 79.8 \\
    2025-06 & 3,681 & 409 & 67.1 & 32.9 & 0.0 & 67.1 & 77.2 \\
    2025-07 & 3,776 & 379 & 71.5 & 28.5 & 0.0 & 71.5 & 80.7 \\
    2025-08 & 3,628 & 348 & 60.1 & 39.9 & 0.0 & 60.1 & 72.9 \\
    2025-09 & 4,320 & 317 & 64.4 & 35.6 & 0.0 & 64.4 & 75.0 \\
    2025-10 & 4,273 & 287 & 61.7 & 38.3 & 0.0 & 61.7 & 78.3 \\
    2025-11 & 3,022 & 256 & 58.5 & 41.5 & 0.0 & 58.5 & 76.2 \\
    2025-12 & 5,493 & 226 & 50.4 & 49.6 & 0.0 & 50.4 & 80.7 \\
    2026-01 & 4,309 & 195 & 59.8 & 40.2 & 0.0 & 59.8 & 73.8 \\
    2026-02 & 4,616 & 164 & 66.0 & 34.0 & 0.0 & 66.0 & 78.2 \\
    2026-03 & 6,234 & 136 & 70.2 & 25.5 & 4.3 & 70.2 & 80.6 \\
    2026-04 & 5,810 & 105 & 67.6 & 27.7 & 4.7 & 67.6 & 81.2 \\
    2026-05 & 6,940 & 75 & 55.2 & 38.1 & 6.7 & 56.7 & 80.7 \\
    2026-06 & 7,947 & 44 & 53.8 & 40.7 & 5.5 & 53.9 & 78.2 \\
    2026-07 & 8,953 & 14 & 33.5 & 36.0 & 30.4 & 33.7 & 76.4 \\
    \hline
  \end{tabular}
  \end{center}
  \footnotesize
  経過日は各月の15日を基準とした概算．Def.は \texttt{Deferred}，NVD\%・VC\%はCPE適用条件の付与率．
\end{table}

\section{追加の限界事項}
\label{app:limits}

\paragraph{VulnCheckの独自エンリッチメントの詳細．}
同社は，約500のCVE関連データソースを正規化した自社基盤をもとにCPEの生成と検証を自動化しており，NISTが公式に登録したCPEのみを用いると述べている\cite{vulncheck-cpe}．
一方で生成規則の詳細や精度指標は公表されていないため，本稿は入出力の観測に基づく評価にとどまる．

\paragraph{\texttt{sourceIdentifier} はCNAの代理指標である．}
表\ref{tab:cna}の分類には，採番CNA以外（ADPや参照元）に由来する識別子が含まれる可能性がある．
厳密なCNA別集計にはCVEレコードのCNAコンテナを直接参照する必要がある．
また表\ref{tab:cna}のApache CNAはhttpd以外のApacheプロジェクト全体を含むため，httpd単体の傾向とは一致しない．

\paragraph{取下げエントリの代替IDを追跡していない．}
\texttt{vulnStatus == "Rejected"} を除外する処理を入れているが，取下げの理由は一様でなく，重複を理由とする取下げでは代替IDへの追跡が必要である．

\paragraph{VulnCheck NVD++ の系列選択．}
本稿が用いたのは \texttt{nist-nvd2} 系列であり，同社が提供する他の系列（\texttt{nist-nvd} など）との差異は評価していない．

\paragraph{CVRFは修正版を持たない．}
CVRF基準では過剰検知（C3b）を測定できない．CSAFは \texttt{known\_not\_affected} を持つが173件に限られる．

\section{Fortinet CVRF基準の年別内訳}
\label{app:cvrf-year}

\begin{table}[htbp]
  \caption{CVRF基準の一致状況（年別，単位\%．各セルは NVD / VulnCheck）}
  \ecaption{Agreement under the CVRF baseline, by year.}
  \label{tab:cvrf-year}
  \begin{center}
  \small
  \begin{tabular}{lrrrrrr}
    \hline
    公開年 & CVE & 組 & 一致 & C2 & C3a & 未検知率 \\
    \hline
    2022-2023 & 254 & 367 & 48.2 / 54.5 & 18.8 / 16.3 & 33.0 / 29.2 & 19.3 / 14.0 \\
    2024 & 117 & 225 & 52.9 / 61.3 & 35.6 / 28.9 & 11.6 / 9.8 & 28.5 / 23.1 \\
    2025 & 204 & 408 & 73.8 / 85.8 & 16.4 / 9.1 & 9.8 / 5.1 & 14.5 / 7.9 \\
    2026 & 83 & 140 & 66.4 / 93.6 & 22.9 / 3.6 & 10.7 / 2.9 & 22.3 / 4.4 \\
    \hline
    合計 & 658 & 1,140 & 60.5 / 71.8 & 21.8 / 14.6 & 17.7 / 13.5 & 19.7 / 12.4 \\
    \hline
  \end{tabular}
  \end{center}
  \footnotesize
  採番元で層別していないため組込み部品のCVEを含む．
  2024年でC2が高いのは，この年のアドバイザリに組込み部品のCVEが多く含まれるためである．
\end{table}


\end{document}
